% SOURCE_AUTHORITY: sga5_fr_workpass.tex, lines 13477--13519 (LF-normalized unit).
% SOURCE_SHA256: 791F4EFFC5E02832D5D77ED03518C8156D6F07E4C8238B03545DB93D883FBB28
A demonstração da proposição 4.6 utiliza essencialmente o seguinte lema:

\medskip
\noindent
\textbf{Lema 4.7.} \emph{Com as hipóteses e notações da proposição 4.6, suponha-se ainda que $\varphi = \id_G$, ou seja, que a ação de $G$ sobre $R$ comuta com $f'_R$, e que $r = \nu(f'_R) \geq 2$.}

\emph{Seja
\[
c = \card(\im(G \longrightarrow \Aut_k(\fm/\fm^2))) \quad .
\]
Tem-se então a seguinte relação de divisibilidade:
\[
c \mid (r - 1)
\]}

\medskip

\begin{prooftext}[Demonstração]
É claro que o inteiro $r$ satisfaz as seguintes condições:
\begin{enumerate}
\item[a)] o homomorfismo
\[
f'_r : A/\fm^r \longrightarrow A/\fm^r
\]
induzido por $f'_R$ é a identidade.
\item[b)] o homomorfismo
\[
f'_{r+1} : A/\fm^{r+1} \longrightarrow A/\fm^{r+1}
\]
é distinto da identidade.
\end{enumerate}

Além disso, é evidente que $G$ opera sobre todos os espaços vetoriais $\fm^i/\fm^{i+1}$ sobre $k$.

Mostraremos primeiro que, sob as hipóteses do lema, existe um único homomorfismo de $k[G]$-módulos
\[
u(f'_R) : \fm/\fm^2 \longrightarrow \fm^r/\fm^{r+1}
\]
tal que se verificam as duas condições seguintes:
\begin{enumerate}
\item[1)] $u(f'_R) \neq 0$
\item[2)] $u(f'_R)(x \bmod \fm^2) = (f'_R(x) - x) \bmod \fm^{r+1}$, qualquer que seja $x \in \fm$.
\end{enumerate}
